% =====================================================================
%  3. The update rule: step = slope \times learning rate
% =====================================================================
\section{The update rule: step = slope $\times$ learning rate}

\begin{frame}{One parameter: the running example, with numbers}
  \small
  Fix the slope at $m = 0.64$ and let GD find the intercept $b$, starting from $b_0 = 0$ with learning rate $\alpha = 0.1$.
  \[
    \frac{d\,\SSR}{db} = -2\sum_{i}\bigl(y_i - b - 0.64\,x_i\bigr),\quad
    \text{step} = \frac{d\,\SSR}{db}\times\alpha,\quad b_{\text{new}} = b_{\text{old}} - \text{step}.
  \]
  \centering\footnotesize
  \begin{tabular}{@{}c c c c c@{}}
    \toprule
    step $k$ & $b_k$ & slope $d\SSR/db$ & step size & $\SSR(b_{k+1})$ \\
    \midrule
    0 & $0$ & $-5.704$ & $-0.5704$ & $0.878$ \\
    1 & $0.5704$ & $-2.282$ & $-0.2282$ & $0.514$ \\
    2 & $0.7986$ & $-0.913$ & $-0.0913$ & $0.456$ \\
    3 & $0.8898$ & $-0.365$ & $-0.0365$ & $0.446$ \\
    \bottomrule
  \end{tabular}

  \raggedright
  \vspace{4pt}
  Each step is about $0.4\times$ the last (start $\SSR = 3.156$), and $b$ settles at \ans{b^\star = 0.951}.\par
  Steps shrink because the slope flattens near the minimum.
\end{frame}

\begin{frame}{Updating several parameters}
  \begin{columns}[T]
    \begin{column}{0.46\textwidth}
      \centering
      \scalebox{0.9}{% Gradient descent loop, steps (i)-(vi)
\begin{tikzpicture}[
  node distance=0.30cm,
  step/.style={draw=#1, line width=0.8pt, rounded corners=4pt, fill=white,
               drop shadow={opacity=0.15, shadow xshift=0.5pt, shadow yshift=-0.7pt},
               text width=3.8cm, minimum height=0.62cm, align=left, font=\scriptsize\color{mInk},
               inner xsep=10pt},
  step/.default=mNavy,
  badge/.style={circle, fill=#1, text=white, font=\tiny\bfseries, inner sep=1pt, minimum size=10pt},
  arr/.style={-{Stealth[length=5pt]}, line width=0.9pt, mGrey}]
  \node[step] (s1) {take the derivative of the loss for each parameter};
  \node[step, below=of s1] (s2) {pick a starting value for each parameter};
  \node[step, below=of s2] (s3) {plug the current values into the derivatives};
  \node[step=mOrange, below=of s3] (s4) {step size $=$ slope $\times$ learning rate};
  \node[step=mOrange, below=of s4] (s5) {new value $=$ old value $-$ step size};
  \node[step=mTeal, below=of s5] (s6) {stop if the step is tiny or max steps reached};
  \foreach \a/\b in {s1/s2, s2/s3, s3/s4, s4/s5, s5/s6} \draw[arr] (\a) -- (\b);
  \draw[arr] (s6.east) -- ++(0.55,0) |- node[pos=0.25, right=3pt, font=\scriptsize, mGrey] {not yet} (s3.east);
  \foreach \n/\lab/\c in {s1/i/mNavy, s2/ii/mNavy, s3/iii/mNavy, s4/iv/mOrange, s5/v/mOrange, s6/vi/mTeal}
    \node[badge=\c] at (\n.west) {\lab};
\end{tikzpicture}
}
    \end{column}
    \begin{column}{0.52\textwidth}
      \small
      For $\hat y = b + mx$ both parameters move:
      \begin{align*}
        \frac{\partial\SSR}{\partial b} &= -2\sum_i\bigl(y_i - b - m x_i\bigr),\\
        \frac{\partial\SSR}{\partial m} &= -2\sum_i x_i\bigl(y_i - b - m x_i\bigr).
      \end{align*}
      In general, for every $w_j$ at once:
      \[
        \ans{w_j \leftarrow w_j - \alpha\,\frac{\partial J}{\partial w_j}}
      \]
      \textbf{Stop} when the step size is near zero (say $<10^{-3}$) or after a maximum number of steps.
    \end{column}
  \end{columns}
\end{frame}

\begin{frame}{Where the update comes from: first-order Taylor}
  \begin{columns}[T]
    \begin{column}{0.52\textwidth}
      \small
      Newton uses a \emph{second}-order Taylor model. Use the \emph{first}-order one, trusted only near $\x_k$
      (constant $J(\x_k)$ dropped):
      \[
        \x_{k+1} = \argmin_{\x}\Bigl\{\underbrace{\nabla J(\x_k)\T(\x - \x_k)}_{\text{straight-line model}}
          + \underbrace{\tfrac{1}{2\alpha}\norm{\x - \x_k}^2}_{\text{stay close}}\Bigr\}.
      \]
      Setting the gradient of the braces to zero gives
      $\nabla J(\x_k) + \tfrac1\alpha(\x - \x_k) = \mathbf 0$, i.e.
      \[
        \ans{\x_{k+1} = \x_k - \alpha\,\nabla J(\x_k)}
      \]
      so the learning rate controls \alert{how far we trust the straight-line model}.
    \end{column}
    \begin{column}{0.46\textwidth}
      \centering
      \footnotesize
      \begin{tabular}{@{}l c c@{}}
        \toprule
        & GD & Newton \\
        \midrule
        Taylor model & 1st order & 2nd order \\
        uses & $\nabla J$ & $\nabla J$, Hessian $H$ \\
        update & $\x - \alpha\nabla J$ & $\x - H^{-1}\nabla J$ \\
        step length & set by $\alpha$ & set by curvature \\
        \bottomrule
      \end{tabular}

      \vspace{6pt}
      \raggedright\footnotesize Newton's method~\cite{chapra} offers another approach; Section 6 compares their costs.
    \end{column}
  \end{columns}
\end{frame}

\begin{frame}{Gradient descent as forward Euler}
  \small
  Let a point slide downhill in continuous time. The \textbf{gradient flow} ODE always lowers $J$:
  \[
    \dot\x(t) = -\nabla J(\x(t)) \quad\Longrightarrow\quad \frac{d}{dt}J(\x(t)) = \nabla J\T\dot\x = -\norm{\nabla J}^2 \le 0 .
  \]
  Solve it with \textbf{forward Euler}, step $h$~\cite{chapra}: $\;\x_{k+1} = \x_k + h\,(-\nabla J(\x_k))$.
  That is \alert{gradient descent with $\alpha = h$}, so Euler's stability analysis applies.

  \vspace{3pt}
  \centering\footnotesize
  \begin{tabular}{@{}l l@{}}
    \toprule
    Numerical ODEs & Optimisation \\
    \midrule
    Euler is stable only if $|1 - h\lambda| < 1$ & the learning-rate limit $\alpha < 2/f''$ (Section 5) \\
    stiff system (eigenvalues far apart) & stretched contours and the zig-zag (Section 4) \\
    implicit (backward) Euler & the proximal-point method, stable for any step \\
    a ball with mass and friction & momentum (Section 8) \\
    $\ddot\x + \tfrac3t\dot\x + \nabla J = 0$ & Nesterov's method~\cite{su2016} \\
    \bottomrule
  \end{tabular}
\end{frame}

\statement{THE ODE CONNECTION}{Gradient descent is forward Euler\\ on the gradient flow.}{So the learning-rate limit is a stability limit, and a stretched valley is a stiff ODE.}

\begin{frame}[fragile]{Many parameters at once: the vectorised update}
  \begin{columns}[T]
    \begin{column}{0.4\textwidth}
      \small
      Multiple linear regression~\cite{chapra}: $\hat y = w_0 + w_1x_1 + \dots + w_nx_n$ with the mean squared error
      \[
        J(\w) = \frac{1}{2m}\sum_{i=1}^{m}\bigl(\hat y_i - y_i\bigr)^2 .
      \]
      Stack a column of ones and the features into $\bX$. All partial derivatives come from one product,
      \[
        \nabla J(\w) = \tfrac1m\,\bX\T(\bX\w - \vy),
      \]
      so one step costs $O(mn)$: two matrix-vector products, no matrix inversion.
    \end{column}
    \begin{column}{0.58\textwidth}
\begin{lstlisting}[style=py]
import numpy as np

def gd(X, y, alpha=0.1, tol=1e-8,
       steps=10_000):
    X = np.column_stack([np.ones(len(y)), X])
    w = np.zeros(X.shape[1])            # (ii)
    for k in range(steps):
        g = X.T @ (X @ w - y) / len(y)  # (iii)
        step = alpha * g                # (iv)
        w = w - step                    # (v)
        if np.linalg.norm(step) < tol:  # (vi)
            break
    return w, k
\end{lstlisting}
      \footnotesize The comments match steps (i)--(vi) of the loop.
    \end{column}
  \end{columns}
\end{frame}
